\newpage 
\section{Roots of Polynomials}

When dealing with different kinds of situations and models, most problems can be interpreted as looking for the 
roots of a non-linear polynomial. For this there are plenty of theory and methods.


\subsection{The Bisection Algorithm}

Given a continuous function \(f: \R \to \R\), with \(a_0, b_0 \in \R\) such that 

\[
    f(a_0) < 0, f(b_0) > 0 \quad \text{or} \quad f(a_0) > 0, f(b_0) < 0
\]

this is \(f(a_0) f(b_0) < 0\), then there must be a root in between (Middle Value Theorem).

We define our iteration as

\[
    x_k = \frac{a_k + b_k}{2}
\]

and apply the following rules:

If \(f(x_k) f(a_k) < 0\) then 

\[
    a_{k + 1} = a_{k}, \quad b_{k + 1} = x_k
\]

Else

\[
    a_{k + 1} = x_k, \quad b_{k + 1} = b_k
\]

then we repeat the process. The termination criteria is the number of iterations set for the computation.

\subsection{The Regula-Falsi Algorithm}

This a variation of the \emph{bisection algorithm} which improves the convergence velocity of the 
algorithm. 

We start as before: 

Given a continuous function \(f: \R \to \R\), with \(a_0, b_0 \in \R\) such that 

\[
    f(a_0) < 0, f(b_0) > 0 \quad \text{or} \quad f(a_0) > 0, f(b_0) < 0
\]

this is \(f(a_0) f(b_0) < 0\) implies there must be a zero in between 

We define our iteration as

\[
    x_k = \frac{a_k + b_k}{2}
\]

Now we use the information of \((a_k, f(a_k))\) and \((b_k, f(b_k))\) to define a line

\[
    s(x) = f(a_k) + \frac{f(b_k) - f(a_k)}{b_k - a_k} (x - a_k)
\]

We set it to \(s(x) = 0 = ...\) which gives us

\[
    s(x_k) = 0 \iff x_k = - \frac{b_k - a_k}{f(b_k) - f(a_k)}f(a_k) = \frac{a_kf(b_k) - b_kf(a_k)}{b_k - a_k}
\]

We used the point at which or secant crosses the x-axis as our new estimation giving us the new iteration rules:

\[
    x_k = \frac{a_kf(b_k) - b_kf(a_k)}{b_k - a_k}
\]

If \(f(x_k)f(a_k) < 0\) then 

\[
    a_{k + 1} a_k, \quad b_{k + 1} = x_k
\]

Else
\[
    a_{k + 1} = x_k, \quad b_{k + 1} = b_k
\]


\subsection{The Secant Algorithm}

This is again a sight modification of the previous algorithm given \(x_{ - 1}, x_{0}\). 

\[
    x_{k + 1} = x_k - \frac{f(x_k)}{\frac{f(x_k) - f(x_{k - 1})}{x_k - x_{k - 1}}}
\]

This time we want to find a root and for that we define different secants which will approximate us 
to that specific point.

\section{The Newton Algorithm}

Takin the Taylor Series of our function \(f\) and only considering the first linear term we get 

\[
    f(y) = f(x_k) + (y - x_k)f'(x_k), \quad y =  x_k - \frac{f(x_k)}{f'(x_k)}
\]

The \emph{Newton algorithm} defines the following iteration:

\[
    x_{k + 1} =  x_k - \frac{f(x_k)}{f'(x_k)}
\]

If we would stop at the quadratic term we would get:

\[
    x_{k + 1} =  x_k - \frac{f'(x_k) \pm \sqrt{(f'(x_k))^2 - 2f(x_k)f''(x_k)}}{f''(x_k)}.
\]

The geometric interpretation is that we are approximating the zero by the tangent at the point 
\((x_k, f(x_k))\).

\subsection{The Newton Algorithm Under Banach}

Given the function \(f(x)\) we define our Newton iteration as 

\[
    \Phi(x) = x - \frac{f(x)}{f'(x)}
\]

with 

\[
    \Phi'(x) = \frac{f(x)f''(x)}{(f'(x))^2} 
\]

\[
    \Phi''(x) = \frac{((f'(x))^3 f''(x) + f(x) f'(x))^2f'''(x) - 2f(x) f'(x)(f''(x))^2}{(f'(x))^4}
\]

We can show that our new function is a contraction by showing: 

\[
    \metric{\Phi'(x)} \le \alpha < 1
\]

For this, we must find an upper bound for \(\metric{f(x)}, \metric{f''(x)}\) or a lower bound for \(\metric{f'(x)}^2\)

\textbf{Example:}

Show that the Newton algorithm is a contraction for \(f(x) = x + \ln(x) - 2\) on \(X = [1, 2]\)

\[
    f'(x) = 1 + \frac{1}{x}, \quad f''(x) = -\frac{1}{x^2}
\]

Returns 

\[
    \Phi'(x) = \frac{f(x)f''(x)}{(f'(x))^2} = \frac{-\frac{1}{x^2}(x + \ln(x) - 2)}{\left(1 + \frac{1}{x}\right)^2} = \frac{-(x + \ln(x) - 2)}{\left(x + 1\right)^2}
\]

Now with \(|\cdot|\)

\[
    |\Phi'(x)| = \frac{|(x + \ln(x) - 2)|}{\left(x + 1\right)^2}
\]

We use the derivatives of \(f(x)\) to analyze the way we are going to find our bounds.

\(f'(x)\) is monotonously increasing, therefore

\[
    f'(x) = 1 + \frac{1}{x} > 0, \quad \text{for} \quad x \in [1, 2] \implies f(x) \in [f(1), f(2)] 
\]

For our denominator we know that it becomes minimal when \(x\) is minimal, therefore

\[
    |\Phi'(x)| = \frac{|f(x)|}{(x + 1)^2} \le \frac{1}{4} = \alpha
\]

\subsection{Convergence Of The Newton ALgorithm}

If \(f: \R \to \R\) is two times continuously differentiable with \(f(x)= 0, f'(x) \ne 0\) i.e. \(x\)  is a root of \(f\). 
If \(x_0 \in [x - \varepsilon, x + \varepsilon]\) where \(\varepsilon\) is infinitesimally small, then the Newton algorithm converges 
\emph{quadratically}.


\subsection{Taylor Expansion For Zero Roots of the Newton ALgorithm}

If we have the case that \(x\) is zero root of \(f(x), f'(x), \dots\) or even leads to un-definitions then we can write \(f\) as 
a Taylor polynomial 

\[
    f(y) = (y - x)^q g(y)
\]

where \(q - 1\) is the grade of last derivate such that \(f'(x) = 0, f''(x) = 0, \dots, f^{(q - 1)} = 0, f^{(q)}\ne 0\)

This gives us 

\[
    f'(x) = q(y - x)^{q - 1} g(y) + (y - x)^q g'(x)
\]

Then 

\[
    \frac{f(y)}{f'(x)} = \frac{(y - x)^q g(y)}{q(y - x)^{q - 1} g(y) + (y - x)^q g'(x)} = \frac{(y - x)g(x)}{qg(x) + (y - x)g'(x)}
\]

Resulting into 

\[
    \lim_{y \to x} \frac{f(y)}{f'(y)} = 0.
\]

